% Presentacion - PFM02 / Sesion 01
% Conversion de fracciones a decimales y viceversa

\pdfminorversion=7
\documentclass[
	spanish,
	aspectratio=1610,
	xcolor={dvipsnames,table}
]{beamer}

\geometry{paperwidth=19.2cm,paperheight=12cm}

\usepackage[utf8]{inputenc}
\usepackage[T1]{fontenc}
\usepackage[spanish,es-nodecimaldot]{babel}
\usepackage{lmodern}
\usepackage{booktabs}
\usepackage{array}
\usepackage{graphicx}
\usepackage{ragged2e}
\usepackage{tikz}
\usetikzlibrary{
	arrows.meta,
	positioning,
	calc,
	decorations.pathreplacing,
	shapes.geometric,
	shapes.callouts,
	fit,
	backgrounds,
	overlay-beamer-styles
}
\usepackage{hyperref}

% Datos del alumno y del programa
\newcommand{\studentname}{Mtro. Martin Jonathan de la Cruz Munoz}
\newcommand{\studentshort}{M. de la Cruz Munoz}
\newcommand{\studentid}{00841}
\newcommand{\studentemail}{martin.delacruz@campus.iiiepe.edu.mx}
\newcommand{\programname}{Maestria en Aprendizaje y Ensenanza de las Matematicas}
\newcommand{\programshort}{MAEM}
\newcommand{\coursecode}{PFM02}
\newcommand{\coursename}{Temas selectos de matematicas I}
\newcommand{\coursegroup}{Grupo A}
\newcommand{\courseperiod}{Marzo--julio 2026}
\newcommand{\institutionname}{IIIEPE}
\newcommand{\institutionlongname}{Instituto de Investigacion, Innovacion y Estudios de Posgrado para la Educacion del Estado de Nuevo Leon}
\newcommand{\institutionplace}{Monterrey, Nuevo Leon}
\newcommand{\iiiepelogo}{img/iiiepe/logo-iiiepe-blanco.png}
\newcommand{\iiiepemark}{img/iiiepe/marca-iiiepe.png}
\newcommand{\iiiepedots}{img/iiiepe/logo-puntos.png}

% Paleta institucional
\definecolor{iiiepeNavy}{HTML}{163B5C}
\definecolor{iiiepeBlue}{HTML}{1F6F9F}
\definecolor{iiiepeAqua}{HTML}{20B7A8}
\definecolor{iiiepeOrange}{HTML}{F28C28}
\definecolor{iiiepePaper}{HTML}{F6F8F7}
\definecolor{iiiepeInk}{HTML}{1E2A32}
\definecolor{iiiepeMuted}{HTML}{667085}

\tikzset{
	paso/.style={
		draw=iiiepeBlue,
		rounded corners=2mm,
		fill=iiiepeBlue!8,
		align=center,
		inner sep=5pt,
		font=\small
	},
	dato/.style={
		draw=iiiepeNavy,
		rounded corners=2mm,
		fill=iiiepePaper,
		align=center,
		inner sep=5pt,
		font=\small\bfseries
	},
	resultado/.style={
		draw=iiiepeAqua!80!black,
		rounded corners=2mm,
		fill=iiiepeAqua!12,
		align=center,
		inner sep=5pt,
		font=\small\bfseries
	},
	verificacion/.style={
		draw=iiiepeOrange,
		rounded corners=2mm,
		fill=iiiepeOrange!12,
		align=center,
		inner sep=5pt,
		font=\scriptsize
	},
	nota/.style={
		draw=iiiepeOrange,
		rounded corners=2mm,
		fill=iiiepeOrange!10,
		align=center,
		inner sep=4pt,
		font=\scriptsize
	},
	flecha/.style={
		-{Latex[length=2.5mm]},
		thick,
		iiiepeBlue
	},
	flechaSecundaria/.style={
		-{Latex[length=2.2mm]},
		thick,
		dashed,
		iiiepeOrange
	},
	resaltar/.style={
		draw=iiiepeOrange,
		rounded corners=1.5mm,
		thick,
		inner sep=2pt
	}
}

\newcommand{\cadenaConversion}[4]{%
\begin{tikzpicture}[node distance=1.25cm]
	\node[dato,visible on=<1->] (a) {#1};
	\node[paso, right=of a,visible on=<2->] (b) {#2};
	\node[resultado, right=of b,visible on=<3->] (c) {#3};
	\node[verificacion, below=0.65cm of b,visible on=<4->] (d) {#4};
	\draw[flecha,visible on=<1->] (a) -- (b);
	\draw[flecha,visible on=<2->] (b) -- (c);
	\draw[flechaSecundaria,visible on=<4->] (d) -- (b);
\end{tikzpicture}%
}

\newcommand{\ejercicioTikzTresPasos}[5]{%
\begin{tikzpicture}[node distance=0.95cm and 1.1cm]
	\node[dato,visible on=<1->] (dato) {#1};
	\node[paso, right=of dato,visible on=<2->] (puno) {#2};
	\node[paso, right=of puno,visible on=<3->] (pdos) {#3};
	\node[resultado, right=of pdos,visible on=<4->] (res) {#4};
	\node[verificacion, below=0.75cm of pdos,visible on=<5->] (ver) {#5};
	\draw[flecha,visible on=<1->] (dato) -- (puno);
	\draw[flecha,visible on=<2->] (puno) -- (pdos);
	\draw[flecha,visible on=<3->] (pdos) -- (res);
	\draw[flechaSecundaria,visible on=<5->] (ver) -- (pdos);
\end{tikzpicture}%
}

\mode<presentation>{
	\usetheme{default}
	\usefonttheme{professionalfonts}
	\setbeamertemplate{navigation symbols}{}
	\setbeamertemplate{blocks}[rounded][shadow=false]
}

\hypersetup{
	colorlinks=true,
	urlcolor=iiiepeBlue,
	linkcolor=iiiepeNavy
}

\setbeamercolor{background canvas}{bg=white}
\setbeamercolor{normal text}{fg=iiiepeInk,bg=white}
\setbeamercolor{structure}{fg=iiiepeBlue}
\setbeamercolor{frametitle}{fg=white,bg=iiiepeNavy}
\setbeamercolor{title}{fg=white,bg=iiiepeNavy}
\setbeamercolor{block title}{fg=white,bg=iiiepeBlue}
\setbeamercolor{block body}{fg=iiiepeInk,bg=iiiepePaper}
\setbeamercolor{alerted text}{fg=iiiepeOrange}
\setbeamerfont{title}{size=\LARGE,series=\bfseries}
\setbeamerfont{frametitle}{size=\large,series=\bfseries}
\setbeamerfont{block title}{series=\bfseries}
\setbeamertemplate{itemize item}{\textcolor{iiiepeAqua}{\large$\blacktriangleright$}}
\setbeamertemplate{itemize subitem}{\textcolor{iiiepeOrange}{\scriptsize$\blacksquare$}}
\setbeamertemplate{enumerate item}{\textcolor{iiiepeBlue}{\insertenumlabel.}}

\newcommand{\localLogo}[1]{%
	\IfFileExists{#1}{\includegraphics[height=0.58cm]{#1}}{\textbf{\institutionname}}%
}

\setbeamertemplate{frametitle}{
	\nointerlineskip
	\begin{beamercolorbox}[wd=\textwidth,ht=1.05cm,dp=0.22cm,leftskip=0.35cm,rightskip=0.25cm]{frametitle}
		\begin{minipage}[c]{0.72\textwidth}
			\insertframetitle
		\end{minipage}
		\hfill
		\raisebox{-0.1cm}{\localLogo{\iiiepelogo}}
	\end{beamercolorbox}
	\vspace{-0.04cm}
	{\color{iiiepeOrange}\rule{\textwidth}{1.4pt}}
}

\setbeamertemplate{footline}{
	\leavevmode
	\hbox{%
		\begin{beamercolorbox}[wd=.34\paperwidth,ht=2.7ex,dp=1ex,leftskip=1em]{author in head/foot}
			\color{white}\studentshort
		\end{beamercolorbox}
		\begin{beamercolorbox}[wd=.38\paperwidth,ht=2.7ex,dp=1ex,center]{title in head/foot}
			\color{white}\programshort\ · \coursecode
		\end{beamercolorbox}
		\begin{beamercolorbox}[wd=.20\paperwidth,ht=2.7ex,dp=1ex,rightskip=1em plus 1fil]{date in head/foot}
			\color{white}Matricula \studentid\hfill\insertframenumber/\inserttotalframenumber
		\end{beamercolorbox}
	}
}
\setbeamercolor{author in head/foot}{bg=iiiepeNavy}
\setbeamercolor{title in head/foot}{bg=iiiepeBlue}
\setbeamercolor{date in head/foot}{bg=iiiepeNavy}

\newcommand{\accentbar}{%
	\begin{tikzpicture}[remember picture,overlay]
		\path[use as bounding box] (0,0) rectangle (0,0);
		\fill[iiiepeAqua] (current page.south west) rectangle ([xshift=0.42\paperwidth,yshift=0.08cm]current page.south west);
		\fill[iiiepeOrange] ([xshift=0.42\paperwidth]current page.south west) rectangle ([xshift=\paperwidth,yshift=0.08cm]current page.south west);
	\end{tikzpicture}
}

\newcommand{\sectionframe}[1]{%
	\begin{frame}[plain]
		\begin{tikzpicture}[remember picture,overlay]
			\fill[iiiepeNavy] (current page.south west) rectangle (current page.north east);
			\IfFileExists{\iiiepedots}{%
				\node[opacity=0.13,anchor=east] at ([xshift=0.8cm]current page.east) {\includegraphics[height=9.4cm]{\iiiepedots}};
			}{}
			\fill[iiiepeAqua] ([xshift=0.85cm,yshift=2.1cm]current page.south west) rectangle ([xshift=1.05cm,yshift=6.3cm]current page.south west);
			\fill[iiiepeOrange] ([xshift=1.15cm,yshift=2.1cm]current page.south west) rectangle ([xshift=1.35cm,yshift=5.35cm]current page.south west);
		\end{tikzpicture}
		\vspace{1.8cm}
		\hspace{1.25cm}
		\begin{minipage}{0.72\paperwidth}
			{\color{white}\Huge\bfseries #1}\\[0.25cm]
			{\color{white!75}Sesion 01 · Conversion de fracciones y decimales}
		\end{minipage}
	\end{frame}
}

% Separadores automaticos desactivados para evitar diapositivas repetitivas.
% \AtBeginSection[]{\sectionframe{\insertsectionhead}}

\title[\coursecode]{Conversion de fracciones a decimales y viceversa}
\subtitle{Sesion 01 · \coursename}
\author[\studentshort]{\texorpdfstring{\studentname\\\footnotesize Matricula: \studentid\\\href{mailto:\studentemail}{\studentemail}}{\studentname, Matricula: \studentid}}
\institute[\institutionname]{\institutionlongname\\\coursecode\ · \coursegroup}
\date[\courseperiod]{\texorpdfstring{\institutionplace\\\courseperiod}{\institutionplace, \courseperiod}}

\begin{document}

\begin{frame}[plain]
	\begin{tikzpicture}[remember picture,overlay]
		\fill[iiiepeNavy] (current page.south west) rectangle ([xshift=0.56\paperwidth]current page.north west);
		\fill[iiiepeBlue] ([xshift=0.56\paperwidth]current page.south west) rectangle (current page.north east);
		\IfFileExists{\iiiepedots}{%
			\node[opacity=0.16,anchor=south east] at ([xshift=0.8cm,yshift=-0.3cm]current page.south east) {\includegraphics[height=8.6cm]{\iiiepedots}};
		}{}
		\node[anchor=north west] at ([xshift=0.55cm,yshift=-0.45cm]current page.north west) {\localLogo{\iiiepelogo}};
		\fill[iiiepeOrange] ([xshift=0.58\paperwidth,yshift=1.20cm]current page.south west) rectangle ([xshift=0.92\paperwidth,yshift=1.34cm]current page.south west);
		\fill[iiiepeAqua] ([xshift=0.58\paperwidth,yshift=0.90cm]current page.south west) rectangle ([xshift=0.82\paperwidth,yshift=1.04cm]current page.south west);
	\end{tikzpicture}

	\vspace{1.75cm}
	\begin{minipage}{0.51\paperwidth}
		{\color{white}\Huge\bfseries Conversion de fracciones\\ a decimales\\ y viceversa}\\[0.25cm]
		{\color{white!86}\large PFM02 · Sesion 01}\\[0.35cm]
		{\color{iiiepeOrange}\rule{0.82\linewidth}{1.3pt}}\\[0.45cm]
		{\color{white}\studentname}\\
		{\color{white!80}\footnotesize Matricula: \studentid}
	\end{minipage}
	\hfill
	\begin{minipage}{0.35\paperwidth}
		\raggedleft
		{\color{white}\Large\bfseries \institutionname}\\[0.18cm]
		{\color{white!82}\footnotesize \programname}\\[0.5cm]
		{\color{white}\small \coursecode\ · \coursename}\\
		{\color{white!78}\small \courseperiod}\\[0.5cm]
		{\color{white!74}\footnotesize \institutionplace}
	\end{minipage}
\end{frame}

\begin{frame}{Contenido}
	\tableofcontents
\end{frame}

\section{Actividad}

\begin{frame}{Ficha de la sesion}
	\begin{columns}[T]
		\begin{column}{0.48\textwidth}
			\begin{block}{Contenido}
				Expresion de fracciones como decimales y de decimales como fracciones.
			\end{block}
			\begin{block}{PDA}
				Usa diversas estrategias al convertir numeros fraccionarios a decimales y viceversa.
			\end{block}
		\end{column}
		\begin{column}{0.48\textwidth}
			\begin{block}{Aprendizaje de la secuencia}
				Convertir fracciones a decimales y viceversa, distinguiendo decimales finitos y periodicos.
			\end{block}
			\begin{block}{Producto}
				Procedimientos claros, ejercicios resueltos, situaciones en contexto, evaluacion y retroalimentacion.
			\end{block}
		\end{column}
	\end{columns}
	\accentbar
\end{frame}

\begin{frame}{Criterios incorporados desde la retroalimentacion}
	\small
	\begin{columns}[T]
		\begin{column}{0.49\textwidth}
			\begin{block}{Estructura de secuencia}
				La actividad se organiza en cinco fases: antecedente, explicacion, practica, evaluacion y retroalimentacion.
			\end{block}
			\begin{block}{Evidencia esperada}
				No basta la respuesta: debe observarse estrategia, procedimiento, simplificacion, verificacion e interpretacion.
			\end{block}
		\end{column}
		\begin{column}{0.49\textwidth}
			\begin{block}{Criterios de evaluacion}
				Las actividades y exposiciones deben mostrar claridad matematica, aplicacion contextual y reflexion sobre errores.
			\end{block}
			\begin{block}{Mejora aplicada}
				Se retoma la experiencia de la Sesion 02 para hacer visible la retroalimentacion y los criterios de revision.
			\end{block}
		\end{column}
	\end{columns}
	\accentbar
\end{frame}

\begin{frame}{Ruta pedagogica de la actividad}
	\begin{columns}[T]
		\begin{column}{0.20\textwidth}
			\begin{block}{1. Recuperar}
				Valor posicional, division y simplificacion.
			\end{block}
		\end{column}
		\begin{column}{0.20\textwidth}
			\begin{block}{2. Clasificar}
				Decimal finito, periodico o fraccion decimal.
			\end{block}
		\end{column}
		\begin{column}{0.20\textwidth}
			\begin{block}{3. Convertir}
				Usar valor posicional, division o ecuaciones.
			\end{block}
		\end{column}
		\begin{column}{0.20\textwidth}
			\begin{block}{4. Verificar}
				Comprobar equivalencia y simplificacion.
			\end{block}
		\end{column}
		\begin{column}{0.20\textwidth}
			\begin{block}{5. Aplicar}
				Interpretar medidas, edades y porcentajes.
			\end{block}
		\end{column}
	\end{columns}
	\vspace{0.25cm}
	\begin{block}{Criterio docente}
		El procedimiento es suficiente cuando el estudiante puede explicar que hizo, por que funciona y como comprueba que no cambio la cantidad.
	\end{block}
	\accentbar
\end{frame}

\begin{frame}{Referentes para explicar procedimientos}
	\small
	\begin{block}{Criterios aplicados}
		\begin{itemize}
			\item \textbf{OpenStax:} una fraccion se interpreta como division; para convertirla a decimal se divide numerador entre denominador.
			\item \textbf{LibreTexts:} en decimales periodicos se multiplica por potencias de $10$ y se resta para cancelar la parte repetida.
			\item \textbf{NCTM:} la fluidez procedimental debe conectar conceptos y procedimientos, usar estrategias y verificar el sentido del resultado.
		\end{itemize}
	\end{block}
	\begin{block}{Aplicacion en esta presentacion}
		Cada ejercicio guiado muestra: anticipacion, procedimiento, resultado y verificacion.
	\end{block}
	{\scriptsize
	Fuentes consultadas:
	\begin{itemize}
		\item \href{https://openstax.org/books/prealgebra-2e/pages/5-3-decimals-and-fractions}{OpenStax, 5.3 Decimals and Fractions}.
		\item \href{https://math.libretexts.org/Bookshelves/Applied_Mathematics/Contemporary_Mathematics_\%28OpenStax\%29/03\%3A_Real_Number_Systems_and_Number_Theory/3.04\%3A__Rational_Numbers/3.4.00\%3A_Exercises}{LibreTexts, Converting Repeating Decimals to a Fraction}.
		\item \href{https://www.nctm.org/Standards-and-Positions/Position-Statements/Procedural-Fluency-in-Mathematics/}{NCTM, Procedural Fluency in Mathematics}.
	\end{itemize}}
	\accentbar
\end{frame}

\begin{frame}{Mapa general de conversion}
	\centering
	\begin{tikzpicture}[
		node distance=1.05cm and 1.25cm,
		box/.style={rectangle, rounded corners=3pt, draw=iiiepeNavy, very thick, fill=iiiepePaper, align=center, minimum width=3.2cm, minimum height=1cm},
		decision/.style={diamond, aspect=2.1, draw=iiiepeBlue, very thick, fill=iiiepeAqua!10, align=center, inner sep=2pt},
		arrow/.style={-{Latex[length=2.8mm]}, very thick, draw=iiiepeBlue}
	]
		\node[box] (frac) {Fraccion\\$\frac{a}{b}$};
		\node[decision, right=of frac] (tipo) {¿denominador\\$10,100,1000$?};
		\node[box, above right=of tipo] (pos) {Valor posicional\\contar ceros};
		\node[box, below right=of tipo] (div) {Division\\$a\div b$};
		\node[box, right=of pos] (finito) {Decimal\\finito};
		\node[box, right=of div] (periodico) {Decimal\\finito o periodico};
		\draw[arrow] (frac) -- (tipo);
		\draw[arrow] (tipo) -- node[above,sloped]{si} (pos);
		\draw[arrow] (tipo) -- node[below,sloped]{no} (div);
		\draw[arrow] (pos) -- (finito);
		\draw[arrow] (div) -- (periodico);
	\end{tikzpicture}
	\vspace{0.25cm}
	\begin{block}{Idea clave}
		Convertir no significa cambiar la cantidad: significa escribir el mismo valor en otra representacion.
	\end{block}
	\accentbar
\end{frame}

\section{Fraccion a decimal}

\begin{frame}{Caso 1: fraccion decimal}
	\begin{columns}[T]
		\begin{column}{0.46\textwidth}
			\begin{block}{Que se hace}
				\begin{enumerate}
					\item Identificar cuantos ceros tiene el denominador.
					\item Tomar el numerador.
					\item Colocar el punto decimal contando esa cantidad de lugares desde la derecha.
				\end{enumerate}
			\end{block}
			\[
				\frac{375}{1000}=0.375
			\]
			\begin{block}{Por que funciona}
				$1000$ representa milesimas; por eso $375$ milesimas se escribe $0.375$.
			\end{block}
		\end{column}
		\begin{column}{0.50\textwidth}
			\centering
			\begin{tikzpicture}[scale=0.9]
		
		% Celdas de valor posicional
		\foreach \x/\lab in {0/U,1/d,2/c,3/m}{
			\draw[fill=iiiepePaper] (\x*1.25,0) rectangle ++(1.25,0.9);
			\node at (\x*1.25+0.625,0.45) {\lab};
		}
		
		% Etiquetas verticales
		\node[font=\scriptsize, rotate=90, anchor=east] at (0.625,-0.18) {unidades};
		\node[font=\scriptsize, rotate=90, anchor=east] at (1.875,-0.18) {décimas};
		\node[font=\scriptsize, rotate=90, anchor=east] at (3.125,-0.18) {centésimas};
		\node[font=\scriptsize, rotate=90, anchor=east] at (4.375,-0.18) {milésimas};
		
		% Número decimal
		\node[font=\large\bfseries, color=iiiepeOrange] at (0.625,1.45) {0.};
		\node[font=\Large\bfseries, color=iiiepeBlue] at (1.875,1.45) {3};
		\node[font=\Large\bfseries, color=iiiepeBlue] at (3.125,1.45) {7};
		\node[font=\Large\bfseries, color=iiiepeBlue] at (4.375,1.45) {5};
		
		% Flecha explicativa
		\draw[-{Latex[length=2.5mm]}, very thick, iiiepeOrange] 
			(4.55,2.15) -- (1.55,2.15);
			
		\node[font=\small, align=center] at (3.05,2.55) 
			{$1000$ tiene\\3 ceros};
		
	\end{tikzpicture}
			\vspace{0.2cm}
			\begin{block}{Error a evitar}
				No mover el punto ``a ojo'': el numero de ceros determina el numero de cifras decimales.
			\end{block}
		\end{column}
	\end{columns}
	\accentbar
\end{frame}

\begin{frame}{Caso 2: fraccion no decimal}
	\begin{block}{Idea de partida}
		Si el denominador no es $10$, $100$, $1000$, etc., se divide el numerador entre el denominador. Antes de operar, se estima la magnitud.
	\end{block}
	\begin{columns}[T]
		\begin{column}{0.47\textwidth}
			\[
			\frac{41}{50}=41\div 50
			\]
			\[
			\begin{array}{rcl}
				41.00\div 50 &=& 0.82\\
				50\times 0.8 &=& 40\\
				41-40 &=& 1\\
				1.00\div 50 &=& 0.02
			\end{array}
			\]
		\end{column}
		\begin{column}{0.49\textwidth}
			\begin{block}{Lectura didactica}
				\begin{enumerate}
					\item $41<50$, entonces el resultado debe ser menor que $1$.
					\item La division produce $0.82$.
					\item Se verifica regresando a fraccion.
				\end{enumerate}
			\end{block}
			\begin{block}{Verificacion}
				\[
					0.82=\frac{82}{100}=\frac{41}{50}
				\]
				La fraccion inicial y el decimal final conservan el mismo valor.
			\end{block}
		\end{column}
	\end{columns}
	\accentbar
\end{frame}

\begin{frame}{Representacion visual: $\frac{3}{4}=0.75$}
	\begin{columns}[T]
		\begin{column}{0.48\textwidth}
			\centering
			\begin{tikzpicture}[scale=1.05]
				\draw[fill=iiiepePaper, thick] (0,0) rectangle (4,1);
				\foreach \x in {1,2,3}{\draw[thick] (\x,0) -- (\x,1);}
				\fill[iiiepeAqua!65] (0,0) rectangle (3,1);
				\node at (2,-0.35) {$\frac{3}{4}$ sombreado};
				\draw[-{Latex[length=2.5mm]}, thick, iiiepeBlue] (0,-1.1) -- (4.25,-1.1);
				\foreach \x/\lab in {0/0,1/{0.25},2/{0.50},3/{0.75},4/1}{
					\draw (\x,-1.22) -- (\x,-0.98);
					\node[font=\scriptsize] at (\x,-1.48) {\lab};
				}
				\fill[iiiepeOrange] (3,-1.1) circle (2.2pt);
			\end{tikzpicture}
		\end{column}
		\begin{column}{0.48\textwidth}
			\begin{block}{Tres formas del mismo valor}
				\[
					\frac{3}{4}=\frac{75}{100}=0.75
				\]
			\end{block}
			\begin{itemize}
				\item La imagen muestra magnitud.
				\item La fraccion muestra partes de un entero.
				\item El decimal usa valor posicional.
			\end{itemize}
		\end{column}
	\end{columns}
	\accentbar
\end{frame}

\begin{frame}{Decimales periodicos}
	\begin{columns}[T]
		\begin{column}{0.47\textwidth}
			\begin{block}{Ejemplo}
				\[
					\frac{8}{11}=8\div 11=0.7272\ldots=0.\overline{72}
				\]
			\end{block}
			\begin{enumerate}
				\item La division no termina.
				\item Los residuos se repiten.
				\item Las cifras repetidas forman el periodo.
			\end{enumerate}
			\begin{block}{Criterio}
				$0.\overline{72}$ no es aproximacion: es una forma exacta de escribir una repeticion infinita.
			\end{block}
		\end{column}
		\begin{column}{0.49\textwidth}
			\centering
			\begin{tikzpicture}[
				node distance=1.15cm,
				circ/.style={circle, draw=iiiepeBlue, very thick, fill=iiiepePaper, minimum size=0.95cm},
				arrow/.style={-{Latex[length=2.5mm]}, thick, draw=iiiepeOrange}
			]
				\node[circ] (r8) {$8$};
				\node[circ, right=of r8] (r3) {$3$};
				\node[circ, below=of r3] (r8b) {$8$};
				\node[align=center, left=of r8b] (txt) {residuos\\se repiten};
				\draw[arrow] (r8) -- node[above] {$7$} (r3);
				\draw[arrow] (r3) -- node[right] {$2$} (r8b);
				\draw[arrow] (r8b) -- (r8);
				\node[font=\small] at ($(r3)!0.5!(r8b)+(1.4,0)$) {$72,72,72,\ldots$};
			\end{tikzpicture}
		\end{column}
	\end{columns}
	\accentbar
\end{frame}

\section{Decimal a fraccion}

\begin{frame}{Decimal finito a fraccion}
	\begin{columns}[T]
		\begin{column}{0.48\textwidth}
			\begin{block}{Procedimiento}
				\begin{enumerate}
					\item Quitar el punto decimal.
					\item Colocar como denominador $1$ seguido de tantos ceros como cifras decimales tenga el numero.
					\item Simplificar.
				\end{enumerate}
			\end{block}
		\end{column}
		\begin{column}{0.48\textwidth}
			\[
				4.325=\frac{4325}{1000}
			\]
			\[
				\frac{4325}{1000}=\frac{865}{200}=\frac{173}{40}
			\]
			\begin{block}{Resultado}
				\[
					4.325=\frac{173}{40}
				\]
			\end{block}
			\begin{block}{Verificacion}
				La fraccion final debe quedar simplificada; si se divide $173\div 40$, se recupera $4.325$.
			\end{block}
		\end{column}
	\end{columns}
	\accentbar
\end{frame}

\begin{frame}{Decimal periodico a fraccion}
	\begin{block}{Ejemplo: $0.\overline{18}$}
		\[
			x=0.181818\ldots
		\]
		\[
			100x=18.181818\ldots
		\]
		\[
			100x-x=18.181818\ldots-0.181818\ldots
		\]
		\[
			99x=18 \qquad x=\frac{18}{99}=\frac{2}{11}
		\]
	\end{block}
	\begin{block}{Por que funciona}
		Al restar, la parte decimal periodica se elimina porque ambas expresiones tienen los mismos decimales infinitos.
	\end{block}
	\accentbar
\end{frame}

\begin{frame}{Decimal periodico mixto}
	\begin{block}{Ejemplo: $0.82\overline{1}$}
		\[
			x=0.821111\ldots
		\]
		\[
			1000x=821.111\ldots \qquad 100x=82.111\ldots
		\]
		\[
			1000x-100x=821.111\ldots-82.111\ldots
		\]
		\[
			900x=739 \qquad x=\frac{739}{900}
		\]
	\end{block}
	\begin{itemize}
		\item Se multiplica por $100$ para dejar el punto antes del periodo.
		\item Se multiplica por $1000$ para dejar el punto despues de la primera repeticion.
	\end{itemize}
	\accentbar
\end{frame}

\begin{frame}{Selector de estrategia}
	\centering
	\begin{tikzpicture}[
		node distance=0.9cm and 1.15cm,
		box/.style={rectangle, rounded corners=3pt, draw=iiiepeNavy, very thick, fill=iiiepePaper, align=center, minimum height=0.9cm, minimum width=3.2cm},
		decision/.style={diamond, aspect=2, draw=iiiepeBlue, very thick, fill=iiiepeAqua!12, align=center, inner sep=1pt},
		arrow/.style={-{Latex[length=2.5mm]}, thick, draw=iiiepeBlue}
	]
		\node[box] (inicio) {¿Que tengo?};
		\node[decision, below left=of inicio] (frac) {fraccion};
		\node[decision, below right=of inicio] (dec) {decimal};
		\node[box, below=of frac] (div) {Dividir\\numerador ÷ denominador};
		\node[box, below=of dec] (den) {Usar denominador\\$10,100,1000$\\o eliminar periodo};
		\node[box, below=1.0cm of $(div)!0.5!(den)$] (ver) {Verificar: estimar, simplificar y revisar sentido};
		\draw[arrow] (inicio) -- (frac);
		\draw[arrow] (inicio) -- (dec);
		\draw[arrow] (frac) -- (div);
		\draw[arrow] (dec) -- (den);
		\draw[arrow] (div) -- (ver);
		\draw[arrow] (den) -- (ver);
	\end{tikzpicture}
	\accentbar
\end{frame}

\section{Ejercicios resueltos}

\begin{frame}{Ejercicio guiado 1: fraccion propia}
	\begin{columns}[T]
		\begin{column}{0.48\textwidth}
			\begin{block}<1->{Ejemplo}
				Convertir $\frac{3}{4}$ a numero decimal.
			\end{block}
			\begin{block}<2->{Anticipacion}
				Como $3<4$, el resultado sera menor que $1$; como $3$ esta cerca de $4$, debe estar cerca de $1$.
			\end{block}
			\onslide<2->{
			\[\frac{3}{4}=3\div 4\]
			}
		\end{column}
		\begin{column}{0.48\textwidth}
			\begin{block}<3->{Procedimiento}
				\onslide<3->{
				\[
				\begin{array}{rcl}
					3.00\div 4 &=& 0.75\\
					30\div 4 &=& 7\quad R2\\
					20\div 4 &=& 5\quad R0
				\end{array}
				\]
				}
			\end{block}
			\begin{block}<4->{Verificacion}
				\onslide<4->{\[0.75=\frac{75}{100}=\frac{3}{4}\]}
			\end{block}
		\end{column}
	\end{columns}
	\vspace{-0.12cm}
	\begin{center}
		\resizebox{0.62\textwidth}{!}{%
			\onslide<3->{\cadenaConversion
				{$\frac{3}{4}$}
				{$3.00\div4$}
				{$0.75$}
				{Verificar: $0.75=\frac{75}{100}=\frac{3}{4}$}}
		}
	\end{center}
	\accentbar
\end{frame}

\begin{frame}{Ejercicio guiado 2: fraccion impropia y mixta}
	\begin{columns}[T]
		\begin{column}{0.48\textwidth}
			\begin{block}<1->{Impropia: $\frac{13}{8}$}
				\onslide<1->{\[13\div 8=1.625\]}
				\onslide<2->{
				\[\begin{array}{rcl}
					13\div 8 &=& 1\quad R5\\
					50\div 8 &=& 6\quad R2\\
					20\div 8 &=& 2\quad R4\\
					40\div 8 &=& 5\quad R0
				\end{array}\]
				}
			\end{block}
		\end{column}
		\begin{column}{0.48\textwidth}
			\begin{block}<3->{Mixta: $6\frac{13}{40}$}
				\onslide<3->{Conservar la parte entera y convertir solo la fraccion:}
				\onslide<3->{\[\frac{13}{40}=13\div 40=0.325\]}
				\onslide<4->{\[6\frac{13}{40}=6+0.325=6.325\]}
			\end{block}
			\begin{block}<5->{Criterio de revision}
				\onslide<5->{Una fraccion impropia o mixta debe producir un decimal mayor que $1$.}
			\end{block}
		\end{column}
	\end{columns}
	\vspace{-0.12cm}
	\begin{center}
		\resizebox{0.78\textwidth}{!}{%
		\begin{tikzpicture}[node distance=0.66cm and 0.90cm]
			\onslide<1->{\node[dato] (i1) {$\frac{13}{8}$};
			\node[paso, right=of i1] (i2) {$13\div8$};
			\node[resultado, right=of i2] (i3) {$1.625$};
			\node[verificacion, right=of i3] (i4) {Mayor que $1$ porque $13>8$.};
			\draw[flecha] (i1) -- (i2);
			\draw[flecha] (i2) -- (i3);
			\draw[flechaSecundaria] (i4) -- (i3);}
			\onslide<3->{\node[dato, below=of i1] (m1) {$6\frac{13}{40}$};
			\node[paso, right=of m1] (m2) {$6+(13\div40)$};
			\node[resultado, right=of m2] (m3) {$6.325$};
			\node[verificacion, right=of m3] (m4) {La parte entera se conserva.};
			\draw[flecha] (m1) -- (m2);
			\draw[flecha] (m2) -- (m3);
			\draw[flechaSecundaria] (m4) -- (m3);}
		\end{tikzpicture}}
	\end{center}
	\accentbar
\end{frame}

\begin{frame}{Ejercicio guiado 3: fraccion periodica}
	\begin{columns}[T]
		\begin{column}{0.46\textwidth}
			\begin{block}<1->{Ejemplo}
				Convertir $\frac{22}{90}$ a decimal.
			\end{block}
			\begin{block}<2->{Primero simplificar}
				\onslide<2->{\[\frac{22}{90}=\frac{11}{45}\]}
			\end{block}
			\begin{block}<3->{Anticipacion}
				\onslide<3->{Como $11<45$, el resultado debe ser menor que $1$.}
			\end{block}
		\end{column}
		\begin{column}{0.50\textwidth}
			\begin{block}<4->{Division}
				\onslide<4->{\[\begin{array}{rcl}
					11.00\div 45 &=& 0.2\ldots\\
					110\div 45 &=& 2\quad R20\\
					200\div 45 &=& 4\quad R20\\
					200\div 45 &=& 4\quad R20
				\end{array}\]}
			\end{block}
			\begin{block}<5->{Resultado}
				\onslide<5->{El residuo $20$ se repite; por eso:}
				\onslide<5->{\[\frac{22}{90}=0.2\overline{4}\]}
			\end{block}
		\end{column}
	\end{columns}
	\vspace{-0.12cm}
	\begin{center}
		\resizebox{0.62\textwidth}{!}{%
		\begin{tikzpicture}[node distance=0.78cm and 0.92cm]
			\onslide<1->{\node[dato] (a) {$\frac{22}{90}$};}
			\onslide<2->{\node[paso, right=of a] (b) {Simplificar\\$\frac{11}{45}$};}
			\onslide<4->{\node[paso, right=of b] (c) {Dividir\\$11\div45$};}
			\onslide<5->{\node[resultado, right=of c] (d) {$0.2\overline{4}$};}
			\onslide<4->{\node[verificacion, below=0.60cm of c] (e) {El residuo $20$ se repite.};}
			\onslide<1->{\draw[flecha] (a) -- (b);}
			\onslide<4->{\draw[flecha] (b) -- (c);}
			\onslide<5->{\draw[flecha] (c) -- (d);}
			\onslide<4->{\draw[flechaSecundaria] (e) -- (c);}
		\end{tikzpicture}}
	\end{center}
	\accentbar
\end{frame}

\begin{frame}{Ejercicio guiado 4: decimal finito}
	\begin{columns}[T]
		\begin{column}{0.48\textwidth}
			\begin{block}<1->{Ejemplo}
				\onslide<1->{Convertir $0.525$ a fraccion.}
			\end{block}
			\begin{block}<2->{Procedimiento}
				\onslide<2->{Tiene tres cifras decimales, por eso el denominador inicial es $1000$:}
				\onslide<2->{\[0.525=\frac{525}{1000}\]}
			\end{block}
		\end{column}
		\begin{column}{0.48\textwidth}
			\begin{block}<3->{Simplificacion}
				\onslide<3->{\[\frac{525}{1000}=\frac{105}{200}=\frac{21}{40}\]}
			\end{block}
			\begin{block}<4->{Verificacion}
				\onslide<4->{\[21\div 40=0.525\]}
				\onslide<4->{La fraccion simplificada conserva el mismo valor.}
			\end{block}
		\end{column}
	\end{columns}
	\begin{center}
		\resizebox{0.90\textwidth}{!}{%
			\onslide<2->{\ejercicioTikzTresPasos
				{$0.525$}
				{3 cifras decimales}
				{$\frac{525}{1000}$}
				{$\frac{21}{40}$}
				{Comprobar: $21\div40=0.525$}}
		}
	\end{center}
	\accentbar
\end{frame}

\begin{frame}{Ejercicio guiado 5: decimal periodico mixto}
	\begin{block}<1->{Ejemplo}
		\onslide<1->{Convertir $5.8\overline{72}=5.8727272\ldots$ a fraccion.}
	\end{block}
	\begin{center}
		\resizebox{0.58\textwidth}{!}{%
		\begin{tikzpicture}[node distance=0.55cm]
			\onslide<2->{\node[dato] (x) {$x=5.8727272\ldots$};}
			\onslide<3->{\node[paso, below=of x] (a) {$1000x=5872.7272\ldots$};}
			\onslide<4->{\node[paso, below=of a] (b) {$10x=58.7272\ldots$};}
			\onslide<5->{\node[resultado, below=of b] (c) {$990x=5814$};}
			\onslide<6->{\node[resultado, right=2.2cm of c] (d) {$x=\frac{323}{55}$};}
			\onslide<3->{\draw[flecha] (x) -- node[right, font=\scriptsize] {mover 3 lugares} (a);}
			\onslide<4->{\draw[flecha] (a) -- node[right, font=\scriptsize] {restar} (b);}
			\onslide<5->{\draw[flecha] (b) -- (c);}
			\onslide<6->{\draw[flecha] (c) -- (d);}
		\end{tikzpicture}}
	\end{center}
	\begin{block}<7->{Verificacion}
		\scriptsize
		\onslide<7->{La resta cancela la parte periodica: $1000x-10x=5814$, por eso $990x=5814$ y $x=\frac{5814}{990}=\frac{323}{55}$. El resultado tiene sentido porque queda entre $5$ y $6$.}
	\end{block}
	\accentbar
\end{frame}

\begin{frame}{Banco de respuestas: fracciones a decimal}
	\small
	\begin{columns}[T]
		\begin{column}{0.49\textwidth}
			\begin{block}{Conversiones finitas}
				\begin{itemize}
					\item $\frac{375}{1000}=0.375$
					\item $\frac{1}{8}=0.125$
					\item $\frac{13}{8}=1.625$
				\end{itemize}
			\end{block}
		\end{column}
		\begin{column}{0.49\textwidth}
			\begin{block}{Conversiones mixtas o periodicas}
				\begin{itemize}
					\item $6\frac{13}{40}=6.325$
					\item $\frac{22}{90}=0.2\overline{4}$
					\item $\frac{1}{3}=0.\overline{3}$
				\end{itemize}
			\end{block}
		\end{column}
	\end{columns}
	\vspace{0.2cm}
	\begin{block}{Criterio}
		Las respuestas sirven para verificar; el aprendizaje se demuestra con procedimiento, estimacion y comprobacion.
	\end{block}
	\accentbar
\end{frame}

\begin{frame}{Banco de respuestas: decimales a fraccion}
	\small
	\begin{columns}[T]
		\begin{column}{0.49\textwidth}
			\begin{block}{Decimales finitos}
				\begin{itemize}
					\item $1.6=\frac{8}{5}$
					\item $0.525=\frac{21}{40}$
					\item $0.025=\frac{1}{40}$
				\end{itemize}
			\end{block}
		\end{column}
		\begin{column}{0.49\textwidth}
			\begin{block}{Decimales periodicos}
				\begin{itemize}
					\item $0.\overline{18}=\frac{2}{11}$
					\item $0.5\overline{3}=\frac{8}{15}$
					\item $5.8\overline{72}=\frac{323}{55}$
				\end{itemize}
			\end{block}
		\end{column}
	\end{columns}
	\vspace{0.2cm}
	\begin{block}{Criterio}
		En decimales finitos se usa denominador potencia de $10$; en periodicos se alinea el periodo y se restan ecuaciones.
	\end{block}
	\accentbar
\end{frame}

\section{Situaciones}

\begin{frame}{Situaciones en contexto: valor posicional}
	\small
	\begin{columns}[T]
		\begin{column}{0.49\textwidth}
			\begin{block}{Situacion 01}
				\begin{itemize}
					\item $9$ de cada $100$: $\frac{9}{100}=0.09$.
					\item La misma lectura aplica para $14$, $32$ y $29$ de cada $100$.
				\end{itemize}
			\end{block}
			\begin{block}{Situacion 04}
				$0.16=\frac{16}{100}=\frac{4}{25}$.
			\end{block}
		\end{column}
		\begin{column}{0.49\textwidth}
			\begin{block}{Situacion 05}
				\begin{itemize}
					\item $3.790=\frac{3790}{1000}$.
					\item Simplificacion: $\frac{3790}{1000}=\frac{379}{100}$.
				\end{itemize}
			\end{block}
			\begin{block}{Criterio}
				El numero de cifras decimales determina el denominador inicial.
			\end{block}
		\end{column}
	\end{columns}
	\accentbar
\end{frame}

\begin{frame}{Situaciones en contexto: division y complemento}
	\small
	\begin{columns}[T]
		\begin{column}{0.49\textwidth}
			\begin{block}{Situaciones 02 y 03}
				\begin{itemize}
					\item Edad: $4\frac{5}{12}=4+5\div12=4.41\overline{6}$.
					\item Falta sembrar: $\frac{18}{18}-\frac{7}{18}=\frac{11}{18}=0.6\overline{1}$.
				\end{itemize}
			\end{block}
		\end{column}
		\begin{column}{0.49\textwidth}
			\begin{block}{Situaciones 06 y 09}
				\begin{itemize}
					\item Llave: $\frac{11}{16}=0.6875$.
					\item Hilo: $\frac{7}{40}=\frac{175}{1000}=0.175$.
				\end{itemize}
			\end{block}
		\end{column}
	\end{columns}
	\vspace{0.2cm}
	\begin{block}{Transicion}
		Las situaciones 07, 08, 10, 11 y 12 se trabajan con representaciones porque dependen de lectura de escala, peso, perimetro o diametro.
	\end{block}
	\accentbar
\end{frame}

\iffalse

\begin{frame}{Situaciones en contexto: valor posicional}
	\scriptsize
	\begin{tabular}{p{0.09\textwidth}p{0.54\textwidth}p{0.28\textwidth}}
		\toprule
		\textbf{Sit.} & \textbf{Dato} & \textbf{Respuesta}\\
		\midrule
		01 & Rezago educativo: $9$, $14$, $32$ y $29$ de cada $100$. & $0.09$, $0.14$, $0.32$, $0.29$.\\
		02 & Edad: $4\frac{5}{12}$ anos. & $4.41\overline{6}$ anos.\\
		03 & Eduardo sembró $7$ de $18$ partes; falta sembrar. & $\frac{11}{18}=0.6\overline{1}$.\\
		04 & IVA: $0.16$. & $\frac{16}{100}=\frac{4}{25}$.\\
		05 & Galon: $3.790$ litros. & $\frac{3790}{1000}=\frac{379}{100}$.\\
		06 & Llave espanola: $\frac{11}{16}$ pulgadas. & $0.6875$ pulgadas.\\
		09 & Hilo especial: $\frac{7}{40}$ mm. & $0.175$ mm.\\
		\bottomrule
	\end{tabular}
	\accentbar
\end{frame}

\begin{frame}{Procedimientos para situaciones seleccionadas}
	\scriptsize
	\begin{columns}[T]
		\begin{column}{0.49\textwidth}
			\begin{block}{Valor posicional}
				\begin{itemize}
					\item $9$ de cada $100$: $\frac{9}{100}=0.09$; se repite la misma lectura para $14$, $32$ y $29$.
					\item IVA: $0.16=\frac{16}{100}=\frac{4}{25}$; se simplifica dividiendo entre $4$.
					\item Galon: $3.790=\frac{3790}{1000}=\frac{379}{100}$; tres cifras decimales implican milésimos.
				\end{itemize}
			\end{block}
		\end{column}
		\begin{column}{0.49\textwidth}
			\begin{block}{Division y complemento}
				\begin{itemize}
					\item Edad: $4\frac{5}{12}=4+5\div12=4.41\overline{6}$.
					\item Falta sembrar: $\frac{18}{18}-\frac{7}{18}=\frac{11}{18}=0.6\overline{1}$.
					\item Llave: $\frac{11}{16}=11\div16=0.6875$; hilo: $\frac{7}{40}=\frac{175}{1000}=0.175$.
				\end{itemize}
			\end{block}
		\end{column}
	\end{columns}
	\vspace{0.1cm}
	\begin{block}{Criterio de revision}
		La respuesta final debe acompanarse de estrategia visible y de una comprobacion rapida del valor obtenido.
	\end{block}
	\accentbar
\end{frame}

\fi

\begin{frame}{Representacion: regla y bascula}
	\begin{columns}[T]
		\begin{column}{0.50\textwidth}
			\begin{block}{Medidas como fraccion decimal}
				Leer la medida en decimal y escribir:
				\[
					2.35=\frac{235}{100}
				\]
			\end{block}
			\centering
			\begin{tikzpicture}[scale=0.82]
				\draw[thick] (0,0) rectangle (6,0.65);
				\foreach \x in {0,...,6}{
					\draw[thick] (\x,0.65) -- (\x,0.25);
					\node[font=\scriptsize] at (\x,-0.18) {\x};
				}
				\foreach \x in {0.1,0.2,...,5.9}{\draw (\x,0.65) -- (\x,0.45);}
				\draw[iiiepeOrange, very thick] (2.35,0.9) -- (2.35,-0.05);
				\draw[flechaSecundaria] (2.35,0.9) -- (2.35,1.08);
				\node[font=\small, color=iiiepeOrange] at (2.35,1.2) {$2.35=\frac{235}{100}$};
				\node[nota, text width=2.1cm] at (4.85,1.35) {Lectura directa\\de la escala};
			\end{tikzpicture}
		\end{column}
		\begin{column}{0.46\textwidth}
			\begin{block}{Peso en bascula}
				Leer la aguja como parte de la unidad; por ejemplo:
				\[
					0.75\text{ kg}=\frac{75}{100}=\frac{3}{4}\text{ kg}
				\]
			\end{block}
			\centering
			\begin{tikzpicture}[scale=0.8]
				\draw[thick, fill=iiiepePaper] (0,0) circle (1.45);
				\foreach \ang/\lab in {210/0,270/{1/2},330/1}{
					\draw[thick] (\ang:1.15) -- (\ang:1.35);
					\node[font=\scriptsize] at (\ang:1.72) {\lab};
				}
				\draw[iiiepeOrange, very thick, -{Latex[length=2.5mm]}] (0,0) -- (305:1.05);
				\node[font=\small] at (0,-2.0) {lectura aproximada};
				\node[paso] at (3.15,0.25) {$0.75$ kg};
				\node[resultado] at (5.0,0.25) {$\frac{3}{4}$ kg};
				\node[verificacion] at (4.05,-0.70) {$\frac{75}{100}=\frac{3}{4}$};
				\draw[flecha] (3.75,0.25) -- (4.45,0.25);
				\draw[flechaSecundaria] (4.05,-0.50) -- (4.05,0.00);
			\end{tikzpicture}
		\end{column}
	\end{columns}
	\accentbar
\end{frame}

\begin{frame}{Representacion: perimetro y diametro}
	\begin{columns}[T]
		\begin{column}{0.50\textwidth}
			\begin{block}{Perimetro}
				Sumar los lados con denominador comun:
				\[
					P=a+b+c+d
				\]
			\end{block}
			\centering
			\begin{tikzpicture}[scale=0.78]
				\coordinate (A) at (0,0);
				\coordinate (B) at (3.1,0.35);
				\coordinate (C) at (2.55,2.1);
				\coordinate (D) at (-0.25,1.6);
				\draw[very thick, fill=iiiepeAqua!10] (A)--(B)--(C)--(D)--cycle;
				\node[below] at ($(A)!0.5!(B)$) {$\frac{3}{4}$};
				\node[right] at ($(B)!0.5!(C)$) {$\frac{5}{8}$};
				\node[above] at ($(C)!0.5!(D)$) {$\frac{1}{2}$};
				\node[left] at ($(D)!0.5!(A)$) {$\frac{7}{16}$};
				\node[nota] at (1.55,2.75) {Sumar todos\\los lados};
				\draw[flechaSecundaria] (1.55,2.50) -- (1.55,1.70);
				\node at (1.35,-0.8) {$P=\frac{12+10+8+7}{16}=\frac{37}{16}$};
			\end{tikzpicture}
		\end{column}
		\begin{column}{0.46\textwidth}
			\begin{block}{Diametro de ruedas}
				Convertir cada diametro fraccionario:
				\[
					\frac{5}{8}=5\div 8=0.625
				\]
			\end{block}
			\centering
			\begin{tikzpicture}[scale=0.82]
				\draw[very thick, fill=iiiepePaper] (0,0) circle (1.25);
				\draw[very thick, iiiepeBlue] (-1.25,0) -- (1.25,0);
				\fill[iiiepeOrange] (-1.25,0) circle (2pt);
				\fill[iiiepeOrange] (1.25,0) circle (2pt);
				\node[above] at (0,0.15) {diametro};
				\node[below] at (0,-1.55) {$d=\frac{5}{8}=0.625$};
				\node[nota] at (3.05,0.85) {Fraccion\\$\rightarrow$ division};
				\draw[flechaSecundaria] (2.35,0.60) -- (1.45,0.15);
			\end{tikzpicture}
		\end{column}
	\end{columns}
	\accentbar
\end{frame}

\begin{frame}{Representacion: perimetro de rectangulo}
	\small
	\begin{columns}[T]
		\begin{column}{0.45\textwidth}
			\begin{block}{Datos}
				\[
					l=\frac{15}{4}=3.75,\qquad a=\frac{10}{8}=1.25
				\]
			\end{block}
			\begin{block}{Operacion}
				Usar la formula $P=2(l+a)$.
			\end{block}
		\end{column}
		\begin{column}{0.52\textwidth}
			\centering
			\begin{tikzpicture}[scale=0.86]
				\draw[very thick, fill=iiiepePaper] (0,0) rectangle (4.4,2.1);
				\node[below] at (2.2,0) {$l=3.75$};
				\node[right] at (4.4,1.05) {$a=1.25$};
				\node[paso] at (7.15,1.45) {$P=2(l+a)$};
				\node[resultado] at (7.15,0.35) {$P=10$};
				\draw[flecha] (4.9,1.05) -- (6.15,1.25);
				\draw[flechaSecundaria] (7.15,1.05) -- (7.15,0.70);
			\end{tikzpicture}
		\end{column}
	\end{columns}
	\vspace{0.1cm}
	\begin{block}{Verificacion}
		$2(3.75+1.25)=2(5)=10$; el resultado queda expresado como numero decimal.
	\end{block}
	\accentbar
\end{frame}

\section{Evaluacion}

\begin{frame}{Autoevaluacion}
	\scriptsize
	\begin{tabular}{p{0.56\textwidth}ccc}
		\toprule
		\textbf{Indicador} & \textbf{3} & \textbf{2} & \textbf{1}\\
		\midrule
		Convierto fracciones a decimales usando el algoritmo aprendido. & $\square$ & $\square$ & $\square$\\[0.2cm]
		Convierto decimales finitos y periodicos a fracciones. & $\square$ & $\square$ & $\square$\\[0.2cm]
		Reconozco y escribo decimales periodicos con la notacion correcta. & $\square$ & $\square$ & $\square$\\[0.2cm]
		Verifico equivalencia, simplificacion y sentido contextual. & $\square$ & $\square$ & $\square$\\
		\bottomrule
	\end{tabular}
	\vspace{0.25cm}
	\begin{block}{Cierre metacognitivo}
		Que estrategia use mejor: valor posicional, division o eliminacion del periodo. Que error debo vigilar en mi siguiente intento.
		\iffalse
		¿Que estrategia use mejor: valor posicional, division o eliminacion del periodo? ¿En que tipo de ejercicio necesito mas practica?
		\fi
	\end{block}
	\accentbar
\end{frame}

\begin{frame}{Retroalimentacion para mejorar}
	\small
	\begin{columns}[T]
		\begin{column}{0.49\textwidth}
			\begin{block}{Errores frecuentes}
				\begin{itemize}
					\item Dividir denominador entre numerador.
					\item Colocar mal el punto decimal.
					\item Omitir la barra del periodo.
					\item No simplificar la fraccion final.
				\end{itemize}
			\end{block}
		\end{column}
		\begin{column}{0.49\textwidth}
			\begin{block}{Preguntas de revision}
				\begin{itemize}
					\item Que estrategia correspondia al tipo de numero?
					\item El resultado conserva la misma cantidad?
					\item La respuesta puede simplificarse?
					\item El resultado tiene sentido en el contexto?
				\end{itemize}
			\end{block}
		\end{column}
	\end{columns}
	\accentbar
\end{frame}

\section{Cierre}

\begin{frame}{Conclusiones}
	\begin{enumerate}
		\item Una fraccion decimal se convierte con valor posicional.
		\item Una fraccion no decimal se convierte dividiendo numerador entre denominador.
		\item Un decimal finito se escribe con denominador potencia de $10$ y se simplifica.
		\item Un decimal periodico se convierte eliminando el periodo mediante resta de ecuaciones.
		\item Los procedimientos se presentan con estrategia, desarrollo y verificacion, siguiendo los criterios de OpenStax, LibreTexts y NCTM.
		\item La mejora incorpora evaluacion y retroalimentacion como parte del aprendizaje, no solo como cierre administrativo.
	\end{enumerate}
	\accentbar
\end{frame}

\begin{frame}[plain]
	\begin{tikzpicture}[remember picture,overlay]
		\fill[iiiepeNavy] (current page.south west) rectangle (current page.north east);
		\IfFileExists{\iiiepedots}{%
			\node[opacity=0.13,anchor=south east] at ([xshift=0.8cm,yshift=-0.2cm]current page.south east) {\includegraphics[height=9cm]{\iiiepedots}};
		}{}
		\node[anchor=north west] at ([xshift=0.55cm,yshift=-0.45cm]current page.north west) {\localLogo{\iiiepelogo}};
	\end{tikzpicture}
	\vspace{2.2cm}
	\begin{center}
		{\color{white}\Huge\bfseries Gracias}\\[0.55cm]
		{\color{white}\Large \studentname}\\[0.18cm]
		{\color{white!78}\footnotesize Matricula \studentid\ · \studentemail}\\[0.45cm]
		{\color{iiiepeOrange}\rule{0.36\paperwidth}{1.4pt}}\\[0.35cm]
		{\color{white!72}\footnotesize \institutionname\ · \coursecode\ · \coursename}
	\end{center}
\end{frame}

\end{document}
